\documentclass[10pt,aspectratio=169]{beamer}
\usepackage{bma01}

\title[Gráficas polares especiales]{Gráficas en coordenadas polares especiales}
\subtitle{Caracoles, rosas, lemniscatas, espirales y cónicas. Ejemplos\\[2pt]\small Semana 15 -- Sesión 1}

\begin{document}

\begin{frame}
  \titlepage
\end{frame}

\begin{frame}{Objetivos de la sesión}
  \begin{enumerate}
    \item Reconocer y graficar las familias clásicas: \resaltar{caracoles}, \resaltar{rosas}, \resaltar{lemniscatas} y \resaltar{espirales}.
    \item Escribir las \resaltar{cónicas} en forma polar y relacionarlas con las órbitas planetarias.
    \item Hallar \resaltar{intersecciones} de curvas polares (incluido el polo).
  \end{enumerate}
  \vspace{1mm}
  {\small\textbf{Competencia 4 del sílabo:} Evalúa la aproximación con polinomios de Taylor y la utilidad de otros sistemas de referencia.}
  \begin{block}{Plan (2 horas; incluye verificación previa de 5 min)}
    \begin{tabular}{ll}
      0:00 -- 0:25 & Caracoles (limaçons) y cardioides\\
      0:25 -- 0:45 & Rosas y lemniscatas\\
      0:45 -- 1:05 & Espirales\\
      1:05 -- 1:30 & Cónicas en polares y órbitas\\
      1:30 -- 1:45 & Intersecciones de curvas polares\\
      1:45 -- 2:00 & Actividad colaborativa, resumen y cierre metacognitivo
    \end{tabular}
  \end{block}
\end{frame}

\begin{frame}{Aula invertida: preparación y verificación}
  \small
  \begin{previa}[antes de la sesión]
    \begin{enumerate}
      \item Lea sobre caracoles, rosas, lemniscatas, espirales y cónicas en polares.
      \item Bosqueje $r=\cos3\theta$ e identifique el número de pétalos.
      \item Pregunta guía: ¿cómo se hallan las intersecciones de dos curvas polares?
    \end{enumerate}
  \end{previa}
  \begin{block}{Verificación al inicio (5 min, en pares)}
    \begin{enumerate}
      \item ¿Cuántos pétalos tiene $r=\cos2\theta$?
      \item ¿Qué simetrías tiene $r^2=\cos2\theta$?
    \end{enumerate}
    Los pares discuten y votan; el docente aclara solo lo que genera dudas y ajusta el ritmo de la sesión.
  \end{block}
\end{frame}

\section{Caracoles}

\begin{frame}{Caracoles: $r=a+b\cos\theta$ \ (ó $a+b\sen\theta$), \ $a,b>0$}
  \centering
  \begin{tikzpicture}
    \begin{axis}[bma, name=c1, xmin=-1.5, xmax=3.5, ymin=-2.2, ymax=2.2, axis equal image, width=4.6cm, height=4.6cm, xtick=\empty, ytick=\empty, xlabel={}, ylabel={}, samples=200, title={\scriptsize $a<b$: lazo interior}]
      \addplot[curva, domain=0:360] ({(1+2*cos(x))*cos(x)},{(1+2*cos(x))*sin(x)});
    \end{axis}
    \begin{axis}[bma, at={(c1.east)}, anchor=west, xshift=0.3cm, name=c2, xmin=-1, xmax=3, ymin=-2, ymax=2, axis equal image, width=4.6cm, height=4.6cm, xtick=\empty, ytick=\empty, xlabel={}, ylabel={}, samples=200, title={\scriptsize $a=b$: cardioide}]
      \addplot[curva2, domain=0:360] ({(1+cos(x))*cos(x)},{(1+cos(x))*sin(x)});
    \end{axis}
    \begin{axis}[bma, at={(c2.east)}, anchor=west, xshift=0.3cm, name=c3, xmin=-1, xmax=3, ymin=-2, ymax=2, axis equal image, width=4.6cm, height=4.6cm, xtick=\empty, ytick=\empty, xlabel={}, ylabel={}, samples=200, title={\scriptsize $b<a<2b$: con hendidura}]
      \addplot[uniVerde, very thick, domain=0:360] ({(1.5+cos(x))*cos(x)/1.3},{(1.5+cos(x))*sin(x)/1.3});
    \end{axis}
    \begin{axis}[bma, at={(c3.east)}, anchor=west, xshift=0.3cm, xmin=-1, xmax=3, ymin=-2, ymax=2, axis equal image, width=4.6cm, height=4.6cm, xtick=\empty, ytick=\empty, xlabel={}, ylabel={}, samples=200, title={\scriptsize $a\ge2b$: convexo}]
      \addplot[uniMorado, very thick, domain=0:360] ({(2+cos(x))*cos(x)/1.5},{(2+cos(x))*sin(x)/1.5});
    \end{axis}
  \end{tikzpicture}

  \vspace{2mm}
  \begin{ejemplo}[1: $r=1+2\cos\theta$]
    \begin{itemize}
      \item Simétrico respecto del eje polar. $r=0$ cuando $\cos\theta=-\frac12$: $\theta=\pm\frac{2\pi}3$ (tangentes en el polo).
      \item Para $\frac{2\pi}3<\theta<\frac{4\pi}3$, $r<0$: se dibuja el \textbf{lazo interior} (en el lado de $x>0$).
      \item $r_{\max}=3$ en $\theta=0$; el lazo interior cruza el eje en $r=-1$ ($\theta=\pi$), es decir en el punto $(1,0)$.
    \end{itemize}
  \end{ejemplo}
\end{frame}

\section{Rosas y lemniscatas}

\begin{frame}{Rosas: $r=a\cos n\theta$ \ ó \ $r=a\sen n\theta$}
  \begin{columns}
    \begin{column}{0.4\textwidth}
      \begin{proposicion}
        \begin{itemize}
          \item $n$ \textbf{impar}: $n$ pétalos (se recorren en $[0,\pi]$).
          \item $n$ \textbf{par}: $2n$ pétalos (en $[0,2\pi]$).
          \item Cada pétalo mide $|a|$; pétalos consecutivos separados por $\frac{2\pi}{n}$ ($n$ impar) ó $\frac\pi n$ ($n$ par).
        \end{itemize}
      \end{proposicion}
      \begin{ejemplo}[2]
        $r=\sen3\theta$: 3 pétalos; $r=0$ en $\theta=0,\frac\pi3,\frac{2\pi}3$ (tangentes en el polo); puntas en $\theta=\frac\pi6,\frac\pi2,\frac{5\pi}6$.
      \end{ejemplo}
    \end{column}
    \begin{column}{0.6\textwidth}
      \centering
      \begin{tikzpicture}
        \begin{axis}[bma, name=r1, xmin=-1.2, xmax=1.2, ymin=-1.2, ymax=1.2, axis equal image, width=4.6cm, height=4.6cm, xtick=\empty, ytick=\empty, xlabel={}, ylabel={}, samples=300, title={\scriptsize $r=\sen3\theta$}]
          \addplot[curva2, domain=0:180] ({sin(3*x)*cos(x)},{sin(3*x)*sin(x)});
        \end{axis}
        \begin{axis}[bma, at={(r1.east)}, anchor=west, xshift=0.3cm, xmin=-1.2, xmax=1.2, ymin=-1.2, ymax=1.2, axis equal image, width=4.6cm, height=4.6cm, xtick=\empty, ytick=\empty, xlabel={}, ylabel={}, samples=400, title={\scriptsize $r=\cos4\theta$}]
          \addplot[curva, domain=0:360] ({cos(4*x)*cos(x)},{cos(4*x)*sin(x)});
        \end{axis}
      \end{tikzpicture}
      \begin{tikzpicture}
        \begin{axis}[bma, name=r3, xmin=-1.2, xmax=1.2, ymin=-1.2, ymax=1.2, axis equal image, width=4.6cm, height=4.6cm, xtick=\empty, ytick=\empty, xlabel={}, ylabel={}, samples=400, title={\scriptsize $r=\cos5\theta$}]
          \addplot[uniVerde, very thick, domain=0:180] ({cos(5*x)*cos(x)},{cos(5*x)*sin(x)});
        \end{axis}
        \begin{axis}[bma, at={(r3.east)}, anchor=west, xshift=0.3cm, xmin=-1.2, xmax=1.2, ymin=-1.2, ymax=1.2, axis equal image, width=4.6cm, height=4.6cm, xtick=\empty, ytick=\empty, xlabel={}, ylabel={}, samples=400, title={\scriptsize $r=\sen2\theta$}]
          \addplot[uniMorado, very thick, domain=0:360] ({sin(2*x)*cos(x)},{sin(2*x)*sin(x)});
        \end{axis}
      \end{tikzpicture}
    \end{column}
  \end{columns}
\end{frame}

\begin{frame}{Lemniscatas: $r^2=a^2\cos2\theta$ \ ó \ $r^2=a^2\sen2\theta$}
  \begin{columns}
    \begin{column}{0.55\textwidth}
      \begin{ejemplo}[3: lemniscata de Bernoulli]
        $r^2=4\cos2\theta$.
        \begin{itemize}
          \item Solo hay puntos donde $\cos2\theta\ge0$: $\theta\in\left[-\frac\pi4,\frac\pi4\right]\cup\left[\frac{3\pi}4,\frac{5\pi}4\right]$.
          \item Simétrica respecto del eje polar, del eje $y$ y del polo ($r\to-r$ no cambia $r^2$).
          \item Tangentes en el polo: $\theta=\pm\frac\pi4$.
          \item Cartesiana: $(x^2+y^2)^2=4(x^2-y^2)$.
        \end{itemize}
      \end{ejemplo}
      \small Es el lugar de los puntos cuyo \emph{producto} de distancias a dos focos $(\pm\sqrt2,0)$ es $2$.
    \end{column}
    \begin{column}{0.45\textwidth}
      \centering
      \begin{tikzpicture}
        \begin{axis}[bma, xmin=-2.4, xmax=2.4, ymin=-1.3, ymax=1.3, axis equal image, width=6.5cm, height=4.4cm, xtick={-2,2}, ytick=\empty, samples=200]
          \addplot[curva, domain=-45:45] ({2*sqrt(max(cos(2*x),0))*cos(x)},{2*sqrt(max(cos(2*x),0))*sin(x)});
          \addplot[curva, domain=135:225] ({2*sqrt(max(cos(2*x),0))*cos(x)},{2*sqrt(max(cos(2*x),0))*sin(x)});
          \addplot[gray, dashed, domain=-1.3:1.3] {x};
          \addplot[gray, dashed, domain=-1.3:1.3] {-x};
          \addplot[only marks, mark=*, uniGranate] coordinates {(1.414,0) (-1.414,0)};
        \end{axis}
      \end{tikzpicture}
    \end{column}
  \end{columns}
\end{frame}

\section{Espirales}

\begin{frame}{Espirales}
  \begin{columns}
    \begin{column}{0.33\textwidth}
      \centering
      \begin{tikzpicture}
        \begin{axis}[bma, xmin=-14, xmax=14, ymin=-14, ymax=14, axis equal image, width=4.8cm, height=4.8cm, xtick=\empty, ytick=\empty, samples=300]
          \addplot[curva, domain=0:720] ({rad(x)*cos(x)},{rad(x)*sin(x)});
        \end{axis}
      \end{tikzpicture}\\
      \small \textbf{Arquímedes} $r=\theta$: espiras equidistantes ($2\pi$). Discos de vinilo, resortes de reloj.
    \end{column}
    \begin{column}{0.33\textwidth}
      \centering
      \begin{tikzpicture}
        \begin{axis}[bma, xmin=-9, xmax=13.5, ymin=-10, ymax=7, axis equal image, width=4.8cm, height=4.8cm, xtick=\empty, ytick=\empty, samples=300]
          \addplot[curva2, domain=0:720] ({exp(rad(x)/5)*cos(x)},{exp(rad(x)/5)*sin(x)});
        \end{axis}
      \end{tikzpicture}\\
      \small \textbf{Logarítmica} $r=e^{b\theta}$: corta a los radios con ángulo constante. Nautilo, galaxias, halcones.
    \end{column}
    \begin{column}{0.33\textwidth}
      \centering
      \begin{tikzpicture}
        \begin{axis}[bma, xmin=-0.6, xmax=1.5, ymin=-0.6, ymax=1.3, axis equal image, width=4.8cm, height=4.8cm, xtick=\empty, ytick={1}, samples=300]
          \addplot[uniVerde, very thick, domain=30:900] ({cos(x)/rad(x)},{sin(x)/rad(x)});
          \addplot[gray, dashed, domain=-0.6:1.5] {1};
        \end{axis}
      \end{tikzpicture}\\
      \small \textbf{Hiperbólica} $r=\frac1\theta$: asíntota $y=1$ (pues $y=\frac{\sen\theta}{\theta}\to1$ si $\theta\to0^+$).
    \end{column}
  \end{columns}
  \begin{ejemplo}[4: ángulo constante de la espiral logarítmica]
    Para $r=e^{b\theta}$, el ángulo $\psi$ entre el radio y la tangente cumple $\tan\psi=\dfrac{r}{dr/d\theta}=\dfrac1b$: constante.
  \end{ejemplo}
\end{frame}

\section{Cónicas en polares}

\begin{frame}{Cónicas en coordenadas polares}
  \begin{columns}
    \begin{column}{0.55\textwidth}
      \begin{teorema}
        Con un foco en el polo y directriz $x=d$, la cónica de excentricidad $e>0$ es
        \[ r=\frac{ed}{1+e\cos\theta}. \]
        \begin{itemize}
          \item $0<e<1$: \textbf{elipse}; \ $e=1$: \textbf{parábola}; \ $e>1$: \textbf{hipérbola}.
        \end{itemize}
      \end{teorema}
      \textbf{Prueba.} Cónica: $\frac{|PF|}{|Pl|}=e$, con $|PF|=r$ y $|Pl|=d-r\cos\theta$. Así $r=e(d-r\cos\theta)$. \ $\blacksquare$
      \begin{ejemplo}[5: órbitas de Kepler]
        Los planetas describen elipses con el Sol en un foco: $r=\frac{a(1-e^2)}{1+e\cos\theta}$. Para la Tierra $e\approx0.017$ (casi circular);
        para el cometa Halley $e\approx0.967$.
      \end{ejemplo}
    \end{column}
    \begin{column}{0.45\textwidth}
      \centering
      \begin{tikzpicture}
        \begin{axis}[bma, xmin=-4.5, xmax=2.5, ymin=-3.3, ymax=3.3, axis equal image, width=6.5cm, height=6.5cm, xtick=\empty, ytick=\empty, samples=300, restrict y to domain=-6:6, restrict x to domain=-8:8, unbounded coords=jump]
          \addplot[curva, domain=0:360] ({1/(1+0.5*cos(x))*cos(x)},{1/(1+0.5*cos(x))*sin(x)});
          \addplot[curva2, domain=-150:150] ({1/(1+cos(x))*cos(x)},{1/(1+cos(x))*sin(x)});
          \addplot[uniVerde, very thick, domain=-110:110] ({1/(1+2*cos(x))*cos(x)},{1/(1+2*cos(x))*sin(x)});
          \addplot[only marks, mark=*, black] coordinates {(0,0)};
          \node[uniAzul] at (axis cs:-1.4,1.3) {\scriptsize $e=\frac12$};
          \node[uniGranate] at (axis cs:-3.3,2.7) {\scriptsize $e=1$};
          \node[uniVerde] at (axis cs:1.3,2.6) {\scriptsize $e=2$};
        \end{axis}
      \end{tikzpicture}

      \small Foco en el polo; $ed=1$.
    \end{column}
  \end{columns}
\end{frame}

\section{Intersecciones}

\begin{frame}{Intersecciones de curvas polares}
  \begin{columns}
    \begin{column}{0.55\textwidth}
      \begin{block}{Cuidado}
        Resolver $f(\theta)=g(\theta)$ \textbf{no basta}: un punto puede estar en ambas curvas con \emph{distintas} representaciones
        (en particular, el \textbf{polo}).
      \end{block}
      \begin{ejemplo}[6]
        $r=1+\cos\theta$ y $r=3\cos\theta$.
        \begin{itemize}
          \item $1+\cos\theta=3\cos\theta\Rightarrow\cos\theta=\frac12$: \ $\theta=\pm\frac\pi3$, $r=\frac32$.
          \item Además ambas pasan por el \textbf{polo}: la cardioide en $\theta=\pi$ y la circunferencia en $\theta=\frac\pi2$.
        \end{itemize}
        Intersecciones: $\left(\frac32,\pm\frac\pi3\right)$ y el polo.
      \end{ejemplo}
      \small Recomendación: siempre graficar ambas curvas.
    \end{column}
    \begin{column}{0.45\textwidth}
      \centering
      \begin{tikzpicture}
        \begin{axis}[bma, xmin=-0.6, xmax=3.3, ymin=-1.7, ymax=1.7, axis equal image, width=6.4cm, height=5.8cm, xtick={1.5,3}, ytick=\empty, samples=200]
          \addplot[curva2, domain=0:360] ({(1+cos(x))*cos(x)},{(1+cos(x))*sin(x)});
          \addplot[curva, domain=0:180] ({3*cos(x)*cos(x)},{3*cos(x)*sin(x)});
          \addplot[only marks, mark=*, black] coordinates {(0.75,1.299) (0.75,-1.299) (0,0)};
        \end{axis}
      \end{tikzpicture}
    \end{column}
  \end{columns}
\end{frame}

\begin{frame}{Actividad colaborativa: Galería polar}
  \small
  \begin{activa}[Aprendizaje basado en proyectos y evaluación entre pares]
    Cada grupo elabora un póster digital de una curva polar especial: ecuación, simetrías, rectas tangentes en el polo, una aplicación (por ejemplo órbitas) y su intersección con otra curva. La clase \textbf{coevalúa} con la rúbrica común.
  \end{activa}
  \begin{itemize}
    \item \textbf{Organización:} equipos de 3 a 4 con roles rotativos (coordinador, redactor, verificador, portavoz); 12--15 minutos.
    \item \textbf{Producto:} Póster digital y coevaluación con rúbrica.
    \item \textbf{Retroalimentación:} el docente recorre los equipos, hace preguntas y no da la solución hasta la puesta en común.
  \end{itemize}
\end{frame}

\section{Ejercicios}

\begin{frame}{Ejercicios propuestos}
  \small
  \begin{ejercicio}
    \begin{enumerate}
      \item Identifique y grafique (simetrías, tangentes en el polo, valores extremos de $r$):
        (a) $r=2-3\sen\theta$, \ (b) $r=3+2\cos\theta$, \ (c) $r=2\cos3\theta$, \ (d) $r^2=9\sen2\theta$, \ (e) $r=e^{-\theta/4}$, $\theta\ge0$.
      \item Halle todas las intersecciones de: (a) $r=\sen\theta$ y $r=\cos\theta$; \ (b) $r=1-\cos\theta$ y $r=\cos\theta$; \ (c) $r^2=\cos2\theta$ y $r=\frac{1}{\sqrt2}$.
      \item Identifique la cónica y halle sus vértices: (a) $r=\dfrac{6}{2+\cos\theta}$, \ (b) $r=\dfrac{4}{1-\sen\theta}$, \ (c) $r=\dfrac{3}{1+2\cos\theta}$.
      \item La órbita de Mercurio tiene $a\approx0.387$ UA y $e\approx0.206$. Halle su distancia mínima (perihelio) y máxima (afelio) al Sol.
      \item Pruebe que en la rosa $r=\cos n\theta$ con $n$ par hay $2n$ pétalos.
      \item Halle la pendiente de la tangente a la espiral $r=\theta$ en $\theta=\pi$ y en $\theta=2\pi$.
    \end{enumerate}
  \end{ejercicio}
\end{frame}

\begin{frame}{Metacognición y uso responsable de la IA}
  \small
  \begin{columns}[t]
    \begin{column}{0.5\textwidth}
      \begin{metacog}[cierre, 3 min]
        \begin{enumerate}
          \metacomun
          \item Mirando todo el curso, ¿qué estrategia de estudio me funcionó y cuál cambiaré para el Examen Final?
        \end{enumerate}
      \end{metacog}
    \end{column}
    \begin{column}{0.5\textwidth}
      \begin{iabox}
        Pida a una IA las intersecciones de $r=1+\cos\theta$ y $r=1$. Verifique el polo y las soluciones obtenidas.\\[2pt]
        \textit{Regla:} primero intente usted; luego use la herramienta para contrastar, nunca para sustituir su razonamiento. Anote el cambio en su bitácora.
      \end{iabox}
    \end{column}
  \end{columns}
\end{frame}

\begin{frame}{Resumen de la sesión}
  \renewcommand{\arraystretch}{1.4}
  \begin{center}
    \begin{tabular}{l|l|l}
      \textbf{Familia} & \textbf{Ecuación} & \textbf{Rasgo}\\ \hline
      Caracol & $r=a+b\cos\theta$ & lazo ($a<b$), cardioide ($a=b$)\\
      Rosa & $r=a\cos n\theta$ & $n$ ó $2n$ pétalos\\
      Lemniscata & $r^2=a^2\cos2\theta$ & figura en ocho\\
      Espirales & $r=a\theta$, $r=e^{b\theta}$, $r=\frac1\theta$ & equidistante, equiangular, asintótica\\
      Cónicas & $r=\frac{ed}{1+e\cos\theta}$ & elipse, parábola, hipérbola
    \end{tabular}
  \end{center}
  \begin{itemize}
    \item Intersecciones: resolver $f=g$ \textbf{y} revisar el polo y otras representaciones.
  \end{itemize}
\end{frame}

\begin{frame}{Cierre del curso}
  \begin{columns}[t]
    \begin{column}{0.5\textwidth}
      \textbf{Lo que construimos:}
      \begin{itemize}
        \item Funciones y modelos.
        \item Límites y continuidad (Bolzano, Weierstrass).
        \item La derivada: reglas, TVM, cadena, implícita, inversas.
        \item Aplicaciones: optimización, linealización, razones relacionadas, gráficas, L'Hospital, Taylor.
        \item Curvas paramétricas y polares.
      \end{itemize}
    \end{column}
    \begin{column}{0.5\textwidth}
      \textbf{Lo que viene (Cálculo Integral):}
      \begin{itemize}
        \item Antiderivadas: ``deshacer'' la derivada.
        \item La integral definida: áreas, volúmenes, longitudes (también de curvas paramétricas y polares).
        \item Series de Taylor: polinomios ``infinitos''.
        \item Ecuaciones diferenciales.
      \end{itemize}
    \end{column}
  \end{columns}
  \vspace{3mm}
  \begin{idea}
    La derivada mide el \textbf{cambio instantáneo}; la integral \textbf{acumula} cambios. Juntas son el lenguaje de las ciencias.
  \end{idea}
  \textbf{Siguiente:} Examen Final (semanas 9 a 14).
\end{frame}


\begin{frame}{Evaluación y trabajo posterior}
  \small
  \begin{evalbox}[Práctica del miércoles de la semana 15: PD9]
    \textbf{Cobertura:} semana 14. \textbf{Naturaleza:} práctica dirigida en equipo; por el feriado del miércoles 9/12 se nivela en la sesión de esta semana. El \textbf{Examen Final} (semana 16) cubre las semanas 9 a 14; el sustitutorio (semana 17), todo el curso.
  \end{evalbox}
  \begin{itemize}
    \item \criterios
    \item \declaracionia
    \item \textbf{Portafolio:} guarde el producto de la actividad, la bitácora de metacognición y los ejercicios propuestos de esta sesión; son evidencia de progreso y base de la retroalimentación.
  \end{itemize}
\end{frame}

\end{document}
